\documentclass[10pt]{article}
\usepackage[margin=20mm]{geometry}
\usepackage{amsmath,amssymb,mathtools,fontspec,unicode-math,xurl,hyperref}
\IfFontExistsTF{DejaVu Serif}{\setmainfont{DejaVu Serif}}{\setmainfont{Latin Modern Roman}}
\setmathfont{latinmodern-math.otf}
\setlength{\parindent}{0pt}\setlength{\parskip}{6pt}\sloppy
\title{Initial scientific review: M25-0034}
\author{}
\date{25 October 2026}
\begin{document}
\maketitle
\section*{Identity, version, and outcome}
Jean Bourgain, Alexander Gamburd, and Peter Sarnak, Strong approximation and Diophantine properties of Markoff triples, DOI 10.1090/jams/1061. This review concerns the supplied 35-page IAS author-hosted manuscript, snapshot 25 October 2026, 21:31:05 UTC. The IAS listing calls it a 2024 work and records a February 2025 posting. These dates do not identify a journal-final version. The PDF SHA-256 is\par
> d24a021751b5d26f1b80883ea62bacff1faa8
> 0781341b7daa0e450fe0e2dd04b.\par
Page numbers below are the printed PDF page numbers, equal to the one-based PDF pages. No author TeX is supplied for this item, and no journal-final comparison was made.\par
Overall judgment: Major for a demonstrated failure in the displayed proof of the middle-game estimate, with two separate local Minor findings. This assessment concerns the derivation as written, not a disproof of the principal strong-approximation or divisibility conclusions. The manuscript explicitly offers alternative methods; that counterevidence is retained below.\par
All 35 pages were read as extracted text. Original images of pages 1, 3, 6–11, 14, 16–29, and 30–32 were inspected. Detailed independent checks concentrated on conic algebra, irreducibility hypotheses, Proposition 14’s matrix argument, the cyclotomic lower bound, and the final elementary sieve. Reading all text does not imply independent verification of every cited theorem or every proof step.\par
\section*{Claims, assumptions, and dependencies}
The main finite-field analysis uses \(x^2+y^2+z^2=xyz\), while the introduction and arithmetic application use \(x^2+y^2+z^2=3xyz\). For \(p\ne3\), multiplication of all coordinates by \(3\) gives the correspondence and conjugates the Vieta moves. Small-prime exclusions must be retained.\par
The main chain is: a connected cage and large-order endgame; a power-saving middle-game estimate used to increase rotation order; a cyclotomic argument preventing very small nonzero components; and the exceptional-prime and integer-sieving arguments. Proposition 14 and Corollary 15 are explicitly part of this chain (pages 18–25). The final sieve further depends on precise transitivity results of Meiri–Puder and an equidistribution input of Mirzakhani.\par
For \(\mathcal M^s\), multiplicities are part of the statement. The proof does not resolve the uniqueness conjecture or automatically establish an unweighted density statement. Theorem 19’s displayed meaning is that, for each fixed \(\nu\), the proportion with at most \(\nu\) distinct prime factors tends to zero. It does not assert the conventional extremal-divisor meaning of a highly composite number.\par
\section*{Problem 1 (Major): invalid subgroup verification in Proposition 14}
\subsection*{Claim and location}
Pages 20–21, equations (71)–(76), seek to verify the nonconcentration hypothesis of Theorem 13 for \[\Phi(t)=
 \begin{pmatrix}
 t&-\sigma t^2-4(1-\sigma)\\1&-t
 \end{pmatrix},\qquad \sigma\ne1.\] The argument asserts that every proper subgroup of \(\mathrm{SL}_2(\mathbf F_p)\) has trivial second derived group and reduces the verification to a double commutator. It also specifies a scalar normalization and a family \(g_t\). Three exact checks show that this verification cannot be accepted as written.\par
\subsection*{Check A: the normalization does not give determinant one}
Write \(\Delta(t)=\det\Phi(t)=(1-\sigma)(4-t^2)\). At nonsingular \(s,t\), the displayed matrix is \[M=\Phi(s)^{-1}\Phi(t)\frac{\Delta(s)}{\Delta(t)}.\] Since a scalar multiplies a two-by-two matrix, \[\det M=\frac{\Delta(t)}{\Delta(s)}
 \left(\frac{\Delta(s)}{\Delta(t)}\right)^2
 =\frac{s^2-4}{t^2-4},\] generally not one. For \(p=101\), \(\sigma=2\), \(s=3\), \(t=4\), both matrices are invertible and the determinant is \(5/12=93\) in \(\mathbf F_{101}\). All conditions of the displayed calculation hold. Scalars may be ignored projectively, but not in a claim of membership in \(\mathrm{SL}_2\). A consistent projective argument or a square-root normalization with its field conditions is needed.\par
\subsection*{Check B: equation (72) has an incorrect entry}
Set \(s=\sigma t+4(1-\sigma)/t\). Direct multiplication gives, away from singular parameters, \[\frac{\operatorname{adj}(\Phi(s))\Phi(t)}
 {(1-\sigma)(4/t-t)}
 =
 \begin{pmatrix}
 \sigma(1+\sigma)t+4\sigma(1-\sigma)/t&
 -\sigma^2t^2+4(1-\sigma)^2\\
 1&0
 \end{pmatrix}.\] Equation (72) instead has \(4(1-\sigma)/t\) in the upper-left entry. Printed minus calculated is \(4(1-\sigma)^2/t\), nonzero for \(p>2\), \(\sigma\ne1\), \(t\ne0\). This cannot be absorbed into a projective scalar because the lower-left entries are both one. Thus the subsequent specialized \(g_\pm\) have not been obtained from the actual \(\Phi\) by the stated calculation. The symbolic multiplication is saved in the evidence.\par
\subsection*{Check C: proper subgroups need not be metabelian}
In \(\mathrm{SL}_2(\mathbf F_{97})\), take \[i=\begin{pmatrix}22&0\\0&75\end{pmatrix},\quad
 j=\begin{pmatrix}0&1\\96&0\end{pmatrix},\quad
 u=\begin{pmatrix}59&60\\59&37\end{pmatrix},\quad
 q=\begin{pmatrix}57&7\\90&40\end{pmatrix}.\] Every determinant is one. Exhaustive multiplication closure gives a subgroup \(H\) of order \(48\), whereas \(|\mathrm{SL}_2(\mathbf F_{97})|=912576\), so \(H\) is proper. Closing all commutators gives \(|H'|=24\) and \(|H''|=8\). In particular, \(H''\) contains the nonscalar matrix \(i\). Its projective second derived group has order four, so the exception is not removed by discarding the center.\par
These examples belong to an algebraic family. For \(p\equiv1\pmod8\), choose \(a^2=-1\), \(b^2=2\), and set \[i=\operatorname{diag}(a,-a),\quad
 j=\begin{pmatrix}0&1\\-1&0\end{pmatrix},\quad
 k=ij,\quad u=(-I+i+j+k)/2,\quad q=(i+j)/b.\] The quaternion relations, cyclic conjugation by \(u\), and the transposition up to signs induced by \(q\) give the binary octahedral group, with normal quaternion subgroup and quotient \(S_3\). Its successive derived groups are the binary tetrahedral and quaternion groups. The construction works for arbitrarily large such primes, so a large-prime qualification does not remove the issue. The accompanying program verifies the specified finite example and a binary tetrahedral example independently; no author code is executed.\par
\subsection*{Counterevidence, dependency, and repair}
The cited Bourgain preprint describes a bicommutator characterization of large proper subgroups (arXiv:1208.4008v1, Section 4). Separating bounded exceptional subgroups and showing that the actual polynomial image cannot lie in them is a plausible part of a repair. It is not a completed replacement for (71)–(76), especially after the corrected matrix computation.\par
Page 18 supplies two alternative routes: an auxiliary-polynomial argument and a Corvaja–Zannier estimate. Conditional on the quoted bound \[20\max\{(l_1l_2)^{1/3},\,l_1l_2/p\},\] summing over \(l_2\le l_1\) dividing \(p^2-1\), for \(p^\epsilon\le l_1\le p^{1/2+\delta_0}\) and fixed \(\delta_0<1/2\), gives \[p^{o(1)}O(l_1^{2/3}+l_1^2/p)=o(l_1).\] Thus this quoted estimate can provide the increasing-order mechanism if all its application hypotheses hold. Its original proof was not independently verified here. The principal theorem is not disproved by the present finding.\par
The Major classification concerns a demonstrated invalid derivation in a method explicitly developed for the central middle-game step, used in Corollary 15 and the main proof. It does not require all alternative arguments to fail or repair to be difficult. A complete repair is not verified.\par
\section*{Problem 2 (Minor): Lemma 12 allows a reducible constant-map case}
Location: pages 16–17, Lemma 12 and its fiber-product proof.\par
The assumptions require \(\alpha_1,\alpha_2,\alpha_3\) not all zero and \(\alpha_3\pm2\sqrt{\alpha_1\alpha_2}\ne\pm2\), but allow \(\alpha_1=\alpha_2=0\). Choose \[p=13,\quad d=e=1,\quad \alpha_1=\alpha_2=0,\quad\alpha_3=1.\] The prime is odd, both positive exponents are less than \(p\), the coefficients are not all zero, and \(1\ne\pm2\). Every printed assumption holds. In the two-dimensional torus the equation becomes \[1=x+x^{-1},\qquad x^2-x+1=(x-4)(x-10).\] It is the union of two distinct nonempty vertical tori, reducible even over \(\mathbf F_{13}\). The proposed map \(g(y)\) is constant and is not the degree-\(2e\) cover invoked in the proof. The same degeneration occurs over the algebraic closure for any \(p>3\).\par
Counterevidence and impact. In the nondegenerate application (59), \(\alpha_1=1/4\) and \(\alpha_2=(B_2-B_1)^2/4\) are both nonzero. This counterexample does not refute that application. The supported issue is local overstatement. Requiring \(\alpha_1\alpha_2\ne0\) excludes this example and matches the branch-cycle proof’s setup, but the full corrected monodromy argument and the coincident-coordinate cases remain to be verified. No complete repaired theorem is certified here.\par
\section*{Problem 3 (Minor): Corollary 21 includes inversion}
Location: page 29, Corollary 21; its application to Corollary 22 is on page 30.\par
The printed assertion is \[|\sigma(U_t)\cap U_t|\ll t^{3/4},\qquad
t\mid p-1,\quad t\le p^{3/4},\quad
\sigma\in\mathrm{PGL}_2(\mathbf F_p),\quad \sigma\ne1.\] Take \(\sigma(x)=x^{-1}\), represented by \(\left(\begin{smallmatrix}0&1\\1&0\end{smallmatrix}\right)\). For odd \(p\) this is a nonidentity projective transformation with nonzero determinant and no pole on \(U_t\subset\mathbf F_p^\times\). Inversion permutes all \(t\) roots, so the left side equals \(t\). For arbitrarily large \(t\), Dirichlet’s theorem supplies primes \(p\equiv1\pmod t\) with \(p>t^{4/3}\). Both restrictions on \(t\) hold, while \(t/t^{3/4}=t^{1/4}\) is unbounded. Hence the uniform estimate is false.\par
Counterevidence. Proposition 20 assumes disjoint divisors for \(S\) and \(T\). The pair \(S(x)=x\), \(T(x)=x^{-1}\) violates that assumption, both supports being \(\{0,\infty\}\). Explicitly carrying that restriction into the corollary would exclude this example; its self-contained printed statement only excludes the identity. On page 30 the relevant map is \((b\eta-1)/(\eta-1)\), whose zero and pole avoid \(0,\infty\) for \(b\ne0,1\). Thus the inversion example does not refute the needed nondegenerate case. The \(b=0\) boundary and exact exponent derivation need separate verification. Merely deleting inversion is not asserted to repair every case. This is a local false auxiliary statement without a demonstrated failure of the core application.\par
\section*{Positive checks and additional limits}
\subsection*{Conic and genus calculations}
For \(x=w+w^{-1}\) and \(\kappa=x^2/(x^2-4)\), substitution of \(y=t+\kappa/t\), \(z=tw+\kappa/(tw)\) gives \[x^2+y^2+z^2-xyz=x^2+\kappa(4-x^2)=0.\] The nondegenerate hypotheses \(x\ne0,\pm2\) are essential. The norm calculation similarly explains the elliptic parametrization.\par
In Lemma 9 the quartic \(P(\xi)=\xi^4+2(B_1+B_2)\xi^2+(B_2-B_1)^2\) has four distinct roots if \(B_1B_2(B_1-B_2)\ne0\) and the characteristic is odd: its two roots in \(\xi^2\) are distinct by discriminant \(16B_1B_2\), and neither is zero. This verifies the genus-one double-cover step, not every later irreducibility application.\par
A local enumeration check is also retained: on page 9, omitting \(\eta=0\) from the quadratic-character sum in the stated \(p\equiv3\pmod4\) case changes the result to \(1-\chi(\xi^2-4)\). At \(p=11,\xi=1\), the sum is \(2\) and the elliptic conic meets six other-coordinate fibers; page 10 prints five. An exhaustive finite check is saved. No failure of the subsequent diameter argument is inferred. A local sign in the parabolic rotation formula is likewise recorded in the evidence, not used to justify the Major judgment.\par
\subsection*{Cyclotomic opening}
If \(a,b,c\in[-2,2]\) satisfy \(a^2+b^2+c^2=abc\), then \[a^2+b^2+c^2=|abc|
\le (|a|^3+|b|^3+|c|^3)/3
\le 2(a^2+b^2+c^2)/3,\] forcing all three to vanish. Thus the cyclotomic integer in (82) is nonzero unless the trace triple is zero. Every conjugate has absolute value at most \(12+8=20\), and \(|N(\eta)|\le20^{\varphi(n)}\le20^n\). A prime ideal containing it has norm no larger; combining with \(n=\operatorname{lcm}(l_1,l_2,l_3)\le\max(l_j)^3\) gives the logarithmic bound in the semisimple cases. Parabolic rotations have order at least \(p\) and require a direct case, rather than being lifted as roots of unity of order \(p\) in that unramified argument.\par
\subsection*{Sieve and limits}
For \(p\equiv1\pmod4\), the zero-coordinate fiber has \(2p-2\) nonzero points. Summing the conic sizes gives \(p^2+3p\) total points, so its proportion is \(2/p+O(p^{-2})\). Products over fixed finite \(Q\) give \(\exp(-2L_Q+O(1))\); summing subsets of at most \(\nu\) primes gives the polynomial factor in (102). The proof correctly fixes \(Q,\nu\) before taking the height limit, and only afterwards makes \(L_Q\) large. An exceptional-prime count with exponent less than one gives a convergent exceptional reciprocal sum by dyadic summation, preserving the divergent good-prime sum in the progression \(1\bmod4\).\par
This verifies the elementary sieve conditional on transitivity and equidistribution. The Meiri–Puder block/sign conventions, passage between Vieta-only and permutation-extended actions, elliptic endgame, full auxiliary-polynomial proof, and exact CKSZ dependency remain explicit tasks. They are neither certified nor promoted to unsupported defects.\par
\section*{Version search and conclusion}
The IAS source page and cited Bourgain preprint were consulted. A bounded title/correction search identified no relevant correction; it does not prove that no erratum exists. No finding is attributed to the journal final without comparison.\par
Judgment | Major for the displayed Proposition 14 route; two local Minor findings.\par
Core impact | Invalid nonconcentration verification; principal theorem not disproved and alternate routes remain.\par
Repair | Correct matrix algebra established; complete replacement proof unverified.\par
Adoption | Focused AI acceptance complete; human adjudication pending; no author adoption requested.\par
Remaining | Journal comparison; subgroup repair; branch/coset/action details; external proofs.\par
\section*{Scope and references}
All scientific analysis and source-page inspection were performed by the main assistant. Independent Python/SymPy scripts and outputs, local enumeration, extraction/render logs, web records, and image ledger accompany coverage/findings/provenance JSON. Author code was not executed.\par
The dispatcher-observed Chat / 6 Pro label and the supplied GPT-6 Astra Pro self-description remain separately attributed; backend equivalence is unverified. The original compiler-absence probe preceded the authorized portable Tectonic follow-up. Its build and page inspections remain preserved as historical evidence. The current report is compiled from this exact TeX with Tectonic 0.17.0 in untrusted, cached-only mode. Exact input, dependency, log, PDF and page-image hashes, together with actual inspection of changed or reflowed pages, are recorded in the accompanying provenance. Original-review actions retain their original attribution; focused AI acceptance is recorded separately. Scientific adjudication remains pending, with no human review or author adoption claimed.\par
9 J. Bourgain, A. Gamburd, and P. Sarnak, supplied IAS author manuscript, 35 pages. https://publications.ias.edu/node/2734. J. Bourgain, A modular Szemeredi–Trotter theorem for hyperbolas, arXiv:1208.4008v1, Proposition 2 and Section 4. https://arxiv.org/abs/1208.4008v1. W. Y. Chen, Nonabelian level structures, Nielsen equivalence, and Markoff triples, Annals of Mathematics 199 (2024), 301–443. Publisher metadata/abstract consulted, not entire proof. https://annals.math.princeton.edu/2024/199-1/p05.\par
\end{document}
